% Part: first-order-logic
% Chapter: syntax-and-semantics
% Section: satisfaction

\documentclass[../../../include/open-logic-section]{subfiles}

\begin{document}

\olfileid{fol}{syn}{sat}

\olsection{Relasi Nyukupi \article{formula} \printtoken{S}{formula}
  ing \article{structure} \printtoken{S}{structure}}

\begin{explain}
Gagasan dhasar kang ngubungake ungkapan kayata term lan
!!{formula}s ing siji pihak karo !!{structure}s ing pihak liyane
yaiku \emph{!!{value}} saka term lan \emph{relasi nyukupi}
!!a{formula}. Sacara ora resmi, !!{value} saka term iku !!{element}
saka !!a{structure}---yen term iku mung konstanta, !!{value}-ne
yaiku obyek kang diwenehake marang konstanta dening !!{structure};
yen term kawangun nganggo !!{function}s, !!{value}-ne dietung saka
!!{value}s konstanta lan fungsi kang diwenehake marang simbol fungsi
ing term kasebut. !!^a{formula} \emph{dicukupi} ing !!a{structure}
yen interpretasi kang diwenehake marang predikat ndadekake
!!{formula} iku bener ing domain !!{structure} kasebut. Gagasan
relasi nyukupi iki ditemtokake kanthi induktif: spesifikasi
!!{structure} langsung nemtokake kapan !!{formula}s atom dicukupi,
dene kapan !!{formula} kompleks dicukupi kita tetepake miturut
konektif utawa kuantor utama lan dicukupi utawa orane
!!{subformula}s kang langsung.

Kasus kuantor ing kene rada ruwet, amarga !!{subformula} kang
langsung saka !!{formula} mawa kuantor nduweni !!{variable} bebas,
lan !!{structure}s ora nemtokake !!{value}s saka !!{variable}s.
Kanggo ngrampungi kangelan iki, kita uga ngenalake \emph{pangaji
variabel} lan nemtokake relasi nyukupi
dudu mung relatif marang !!a{structure}, nanging marang
!!a{structure} bebarengan karo pangaji !!a{variable}.
\end{explain}

\begin{defn}[Pangaji Variabel]
\emph{Pangaji variabel}~$s$ kanggo !!a{structure}~$\Struct{M}$
iku fungsi kang masangake saben !!{variable} karo
!!a{element} saka~$\Domain M$, yaiku
$s\colon \Var \to \Domain M$.
\end{defn}

\begin{explain}
!!^a{structure} menehi !!a{value} marang saben !!{constant},
lan pangaji variabel menehi nilai marang saben variabel.
Nanging, kita uga pengin nggunakake term kang kawangun saka
simbol-simbol kasebut kanggo menehi jeneng marang !!{element}s
saka !!{domain}. Mulane, kita nemtokake !!{value} saka term
kanthi induktif. Kanggo !!{constant}s lan variabel, nilai kasebut
persis kaya kang ditemtokake dening !!{structure} utawa pangaji
variabel; kanggo term kang luwih rumit, nilaine dietung kanthi
rekursif nggunakake fungsi kang diwenehake dening !!{structure}
marang !!{function}s.
\end{explain}

\begin{defn}[!!^{value} saka Term]
Yen $t$ iku term saka basa~$\Lang L$, $\Struct M$ iku
!!{structure} kanggo~$\Lang L$, lan $s$ iku pangaji !!a{variable}
kanggo~$\Struct M$, \emph{!!{value}}~$\Value{t}{M}[s]$
ditetepake kaya mangkene:
\begin{enumerate}
\item \indcase{t}{c}{$\Value{\indfrm}{M}[s] = \Assign{\indcomplex}{M}$.}
\item \indcase{t}{x}{$\Value{\indfrm}{M}[s] = s(\indcomplex)$.}
\item \indcase{t}{\Atom{f}{t_1, \ldots, t_n}}{
\[
\Value{\indfrm}{M}[s] = \Assign{f}{M}(\Value{t_1}{M}[s], \ldots,
\Value{t_n}{M}[s]).
\]}
\end{enumerate}
\end{defn}

\begin{defn}[Varian-$x$]
Yen $s$ iku pangaji !!a{variable} kanggo !!a{structure}~$\Struct M$,
mula sembarang pangaji !!{variable}~$s'$ kanggo~$\Struct M$
kang beda saka~$s$ paling akeh mung ing nilai kang diwenehake
marang~$x$ diarani \emph{varian-$x$} saka~$s$. Yen $s'$ iku
varian-$x$ saka~$s$, kita nulis $\varAssign{s'}{s}{x}$.
\end{defn}

\begin{explain}
Gatekna yen varian-$x$ saka pangaji~$s$ ora \emph{kudu}
menehi nilai kang beda marang~$x$. Nyatane, saben pangaji
iku varian-$x$ saka awake dhewe.
\end{explain}

\begin{defn}
  Yen $s$ iku pangaji !!a{variable} kanggo !!a{structure}~$\Struct M$
  lan $m \in \Domain{M}$, mula pangaji~$\Subst{s}{m}{x}$
  yaiku pangaji variabel kang ditetepake dening
  \[\Subst{s}{m}{x}(y) = \begin{cases}
    m & \text{yen } y \ident x\\
    s(y) & \text{yen ora}.
  \end{cases}\]
\end{defn}

Kanthi tembung liya, $\Subst{s}{m}{x}$ iku varian-$x$
tartamtu saka~$s$ kang menehi !!{element} domain~$m$ marang~$x$,
lan menehi nilai marang !!{variable}s saliyane~$x$
padha karo kang diwenehake~$s$.

\begin{defn}[Relasi Nyukupi]
\ollabel{defn:satisfaction}
Relasi nyukupi tumrap !!a{formula}~$!A$ ing !!a{structure}~$\Struct M$
relatif marang pangaji !!a{variable}~$s$, kang ditulis
$\Sat{M}{!A}[s]$, ditetepake kanthi rekursif kaya mangkene.
(Kita nulis $\Sat/{M}{!A}[s]$ kanggo teges "ora $\Sat{M}{!A}[s]$".)
\begin{enumerate}
\tagitem{prvFalse}{%
  \indcase{!A}{\lfalse}{$\Sat/{M}{\indfrm}[s]$.}}{}

\tagitem{prvTrue}{%
  \indcase{!A}{\ltrue}{$\Sat{M}{\indfrm}[s]$.}}{}

\item \indcase{!A}{\Atom{R}{t_1, \dots, t_n}}{$\Sat{M}{\indfrm}[s]$
  yen lan mung yen $\langle \Value{t_1}{M}[s], \dots, \Value{t_n}{M}[s] \rangle \in
  \Assign{R}{M}$.}

\item \indcase{!A}{\eq[t_1][t_2]}{$\Sat{M}{\indfrm}[s]$ yen lan mung yen
  $\Value{t_1}{M}[s] = \Value{t_2}{M}[s]$.}

\tagitem{prvNot}{%
  \indcase{!A}{\lnot !B}{$\Sat{M}{\indfrm}[s]$ yen lan mung yen
    $\Sat/{M}{!B}[s]$.}}{}

\tagitem{prvAnd}{%
  \indcase{!A}{(!B \land !C)}{$\Sat{M}{\indfrm}[s]$ yen lan mung yen
    $\Sat{M}{!B}[s]$ lan $\Sat{M}{!C}[s]$.}}{}

\tagitem{prvOr}{%
  \indcase{!A}{(!B \lor !C)}{$\Sat{M}{\indfrm}[s]$ yen lan mung yen
    $\Sat{M}{!B}[s]$ utawa $\Sat{M}{!C}[s]$ (utawa kalorone).}}{}

\tagitem{prvIf}{%
  \indcase{!A}{(!B \lif !C)}{$\Sat{M}{\indfrm}[s]$ yen lan mung yen
    $\Sat/{M}{!B}[s]$ utawa $\Sat{M}{!C}[s]$ (utawa kalorone).}}{}

\tagitem{prvIff}{%
  \indcase{!A}{(!B \liff !C)}{$\Sat{M}{\indfrm}[s]$ yen lan mung yen
    $\Sat{M}{!B}[s]$ lan $\Sat{M}{!C}[s]$ kalorone,
    utawa $\Sat{M}{!B}[s]$ lan $\Sat{M}{!C}[s]$ kalorone ora kelakon.}}{}

\tagitem{prvAll}{%
  \indcase{!A}{\lforall[x][!B]}{$\Sat{M}{\indfrm}[s]$ yen lan mung yen
    kanggo saben !!{element}~$m \in \Domain M$,
    $\Sat{M}{!B}[\Subst{s}{m}{x}]$.}}{}

\tagitem{prvEx}{%
  \indcase{!A}{\lexists[x][!B]}{$\Sat{M}{\indfrm}[s]$ yen lan mung yen
    kanggo paling ora siji !!{element}~$m \in \Domain M$,
    $\Sat{M}{!B}[\Subst{s}{m}{x}]$.}}{}
\end{enumerate}
\end{defn}

\begin{explain}
Pangaji variabel wigati ing
\iftag{notprvEx,notprvAll}{klausa pungkasan}{rong klausa pungkasan}.
\iftag{prvAll}{ Kita ora bisa nemtokake kapan
$\lforall[x][!B(x)]$ dicukupi kanthi kandha "kanggo kabeh
$m \in \Domain{M}$, $\Sat{M}{!B(m)}$".}{}\iftag{prvEx}{
Kita ora bisa nemtokake kapan $\lexists[x][!B(x)]$ dicukupi kanthi
kandha "kanggo paling ora siji $m \in \Domain{M}$,
$\Sat{M}{!B(m)}$".}{} Sebabe, yen $m \in \Domain M$,
unsur kasebut dudu simbol basa, mula $!B(m)$~dudu !!a{formula}
(yaiku, $\Subst{!B}{m}{x}$ ora ditetepake). Kita uga ora bisa
nganggep yen kita nduweni !!{constant}s utawa term kang bisa menehi
jeneng marang saben !!{element} saka~$\Struct{M}$, amarga definisi
!!{structure}s ora mbutuhake kuwi. Ing basa standar, himpunan
!!{constant}s iku !!{denumerable}, mula yen $\Domain{M}$
ora !!{enumerable}, cacah !!{constant}s malah ora cukup kanggo
menehi jeneng marang saben obyek.

Kita ngrampungi masalah iki kanthi ngenalake pangaji !!{variable},
kang ngidini kita ngubungake variabel langsung karo !!{element}s
domain. Dadi, tinimbang kandha, upamane,
\iftag{prvEx}{$\lexists[x][!B(x)]$}{$\lforall[x][!B(x)]$}
dicukupi ing~$\Struct M$ yen lan mung yen
\iftag{prvEx}{kanggo paling ora siji $m \in \Domain{M}$,
$\Sat{M}{!B(m)}$}{kanggo kabeh $m \in \Domain{M}$,
$\Sat{M}{!B(m)}$}, kita kandha yen formula kasebut dicukupi
ing~$\Struct M$ \emph{relatif marang}~$s$ yen lan mung yen
$!B(x)$ dicukupi relatif marang~$\Subst{s}{m}{x}$
\iftag{prvEx}{kanggo paling ora siji}{kanggo saben}
$m \in \Domain M$. Cathetan koreksi sumber OLPL-073:
cabang prvEx ing ukara bandhingan sumber mung nyebut
anane unsur domain nanging ora nyebut syarat relasi nyukupi
tumrap matriks; syarat kasebut dibalekake supaya sejajar
karo cabang universal lan ukara sadurunge.
\end{explain}

\begin{ex}
Upamakna $\Lang{L} = \{a, b, f, R\}$ kanthi $a$ lan $b$
minangka !!{constant}s, $f$~!!{function} rong panggonan,
lan $R$~!!{predicate} rong panggonan. Gatekna
!!{structure}~$\Struct{M}$ kang ditetepake dening:
\begin{enumerate}
\item $\Domain M = \{1, 2, 3, 4\}$
\item $\Assign{a}{M} = 1$
\item $\Assign{b}{M} = 2$
\item $\Assign{f}{M}(x, y) = x+y$ yen $x+y \le 3$ lan $= 3$
  yen ora.
\item $\Assign{R}{M} = \{\tuple{1, 1}, \tuple{1, 2}, \tuple{2, 3}, \tuple{2, 4}\}$
\end{enumerate}
Fungsi $s(x) = 1$ kang menehi $1 \in \Domain{M}$
marang saben !!{variable} iku pangaji variabel kanggo~$\Struct{M}$.

Banjur
\begin{align*}
\Value{f(a,b)}{M}[s] & = \Assign{f}{M}(\Value{a}{M}[s], \Value{b}{M}[s]).
\intertext{Amarga $a$ lan $b$ iku !!{constant}s, $\Value{a}{M}[s]
  = \Assign{a}{M} = 1$ lan $\Value{b}{M}[s] = \Assign{b}{M} = 2$. Mula}
\Value{f(a,b)}{M}[s] & = \Assign{f}{M}(1, 2) = 1+2 = 3.
\intertext{Kanggo ngetung nilai $f(f(a,b),a)$, kita kudu nimbang}
\Value{f(f(a,b),a)}{M}[s] & = \Assign{f}{M}(\Value{f(a, b)}{M}[s],
\Value{a}{M}[s]) = \Assign{f}{M}(3, 1) = 3,
\intertext{amarga $3+1 > 3$. Amarga $s(x) = 1$ lan $\Value{x}{M}[s] =
  s(x)$, kita uga nduweni}
\Value{f(f(a,b),x)}{M}[s] & = \Assign{f}{M}(\Value{f(a, b)}{M}[s],
\Value{x}{M}[s]) = \Assign{f}{M}(3, 1) = 3,
\end{align*}

!!{formula} atom~$R(t_1, t_2)$ dicukupi yen pasangan nilai
argumene, yaiku $\tuple{\Value{t_1}{M}[s],
  \Value{t_2}{M}[s]}$, dadi !!a{element} saka~$\Assign{R}{M}$.
Dadi, upamane, kita nduweni $\Sat{M}{R(b,f(a,b))}[s]$
amarga $\tuple{\Value{b}{M},
  \Value{f(a,b)}{M}} = \tuple{2, 3} \in \Assign{R}{M}$,
nanging $\Sat/{M}{R(x, f(a,b))}[s]$ amarga
$\tuple{1, 3} \notin \Assign{R}{M}$.
Cathetan koreksi sumber OLPL-072: interpretasi predikat
iku ditemtokake dening struktur, dudu pangaji variabel;
indeks pangaji kang
kelebihan ing pungkasan rumus sumber dicopot.

Kanggo nemtokake apa formula nonatomik~$!A$ dicukupi,
gunakna klausa ing definisi induktif kang cocog karo
konektif utama. Upamane, konektif utama ing
$R(a, a) \lif (R(b, x) \lor R(x, b))$ yaiku~$\lif$, lan
\begin{align*}
  & \Sat{M}{R(a, a) \lif (R(b, x) \lor R(x, b))}[s]
  \text{ yen lan mung yen }\\
  & \qquad
  \Sat/{M}{R(a,a)}[s] \text{ utawa } \Sat{M}{R(b, x) \lor R(x, b)}[s]
  \intertext{Amarga $\Sat{M}{R(a,a)}[s]$ (awit $\tuple{1,1} \in
    \Assign{R}{M}$), kita durung bisa nemtokake wangsulane lan kudu
    mriksa dhisik apa $\Sat{M}{R(b, x) \lor R(x, b)}[s]$:}
  & \Sat{M}{R(b, x) \lor R(x, b)}[s] \text{ yen lan mung yen }\\
  & \qquad \Sat{M}{R(b,x)}[s] \text{ utawa } \Sat{M}{R(x, b)}[s]
  \intertext{Iki pancen bener, amarga $\Sat{M}{R(x, b)}[s]$
    (awit $\tuple{1,2} \in \Assign{R}{M}$).}
\end{align*}

Elinga yen varian-$x$ saka~$s$ iku pangaji variabel kang
beda saka $s$ paling akeh mung ing nilai kang diwenehake
marang~$x$. Kanggo saben !!{element} ing~$\Domain{M}$,
ana varian-$x$ saka~$s$:
\begin{align*}
  s_1 & = \Subst{s}{1}{x}, &
  s_2 & = \Subst{s}{2}{x},\\
  s_3 & = \Subst{s}{3}{x}, &
  s_4 & = \Subst{s}{4}{x}.
\end{align*}
Dadi, upamane, $s_2(x) = 2$ lan $s_2(y) = s(y) = 1$
kanggo kabeh variabel~$y$ saliyane~$x$. Iki kabeh
varian-$x$ saka~$s$ kanggo struktur~$\Struct{M}$,
amarga $\Domain{M} = \{1, 2, 3, 4\}$. Gatekna,
utamane, yen $s_1 = s$ ($s$~mesthi dadi varian-$x$
saka awake dhewe).

\iftag{prvEx}{Kanggo nemtokake apa !!{formula} mawa kuantor
  eksistensial~$\lexists[x][!A(x)]$ dicukupi, kita kudu nemtokake
  apa $\Sat{M}{!A(x)}[\Subst{s}{m}{x}]$ kanggo paling ora siji
  $m \in \Domain M$. Dadi,
  \[
  \Sat{M}{\lexists[x][(R(b,x) \lor R(x,b))]}[s],
  \]
  amarga $\Sat{M}{R(b,x) \lor R(x, b)}[\Subst{s}{1}{x}]$
  ($\Subst{s}{3}{x}$ uga bisa dadi paseksi). Nanging,
  \[
  \Sat/{M}{\lexists[x][(R(b,x) \land R(x,b))]}[s]
  \]
  amarga kanggo sembarang $m \in \Domain{M}$ kang kita pilih,
  $\Sat/{M}{R(b,x) \land R(x,b)}[\Subst{s}{m}{x}]$.}{}

\iftag{prvAll}{Kanggo nemtokake apa !!{formula} mawa kuantor
  universal~$\lforall[x][!A(x)]$ dicukupi, kita kudu nemtokake
  apa $\Sat{M}{!A(x)}[\Subst{s}{m}{x}]$ kanggo kabeh
  $m \in \Domain M$. Dadi,
  \[
  \Sat{M}{\lforall[x][(R(x,a) \lif R(a,x))]}[s],
  \]
  amarga $\Sat{M}{R(x,a) \lif R(a,x)}[\Subst{s}{m}{x}]$
  kanggo kabeh $m \in \Domain M$. Kanggo $m = 1$,
  kita nduweni $\Sat{M}{R(a,x)}[\Subst{s}{1}{x}]$,
  mula konsekuene bener; kanggo $m = 2$, $3$, lan~$4$,
  kita nduweni $\Sat/{M}{R(x,a)}[\Subst{s}{m}{x}]$,
  mula antesedene salah. Nanging,
  \[
  \Sat/{M}{\lforall[x][(R(a,x) \lif R(x,a))]}[s]
  \]
  amarga $\Sat/{M}{R(a,x) \lif R(x,a)}[\Subst{s}{2}{x}]$
  (awit $\Sat{M}{R(a, x)}[\Subst{s}{2}{x}]$ lan
  $\Sat/{M}{R(x, a)}[\Subst{s}{2}{x}]$).}{}

\iftag{defEx}{Kanggo nemtokake apa !!{formula} mawa kuantor
  eksistensial~$\lexists[x][!A(x)]$ dicukupi, kita kudu nemtokake
  apa $\Sat{M}{\lnot \lforall[x][\lnot !A(x)]}[s]$.
  Upamane, kita nduweni
  \[
  \Sat{M}{\lexists[x][(R(b,x) \lor R(x,b))]}[s].
  \]
  Kaping pisan, $\Sat{M}{R(b,x) \lor R(x, b)}[\Subst{s}{1}{x}]$
  ($\Subst{s}{3}{x}$ uga bisa dadi paseksi). Mula,
  $\Sat/{M}{\lnot(R(b,x) \lor R(x,b))}[\Subst{s}{1}{x}]$,
  mula $\Sat/{M}{\lforall[x][\lnot(R(b,x) \lor R(x,b))]}[s]$,
  lan mula $\Sat{M}{\lnot\lforall[x][\lnot(R(b,x) \lor
  R(x,b))]}[s]$. Cathetan koreksi sumber OLPL-077:
  kurung bukak kang keluwihan ing rong rumus kuantor
  dicopot supaya matriks negasine imbang. Ing sisih liya,
  \[
  \Sat/{M}{\lexists[x][(R(b,x) \land R(x,b))]}[s].
  \]
  Cathetan koreksi sumber OLPL-074: koma kang kleru mlebu
  ing argumen formula ing tampilan sumber wis dicopot.
  Iki amarga $\Sat{M}{\lforall[x][\lnot(R(b,x) \land
  R(x,b))]}[s]$, awit ora ana $m \in \Domain M$ kang
  ndadekake $\Sat{M}{R(b,x) \land R(x,b)}[\Subst{s}{m}{x}]$.
  Kaya kang bisa dikira saka tuladha iki,
  $\Sat{M}{\lexists[x][!A(x)]}[s]$ yen lan mung yen
  $\Sat{M}{!A(x)}[\Subst{s}{m}{x}]$ kanggo paling ora
  siji $m \in \Domain M$.}{}

\iftag{defAll}{Kanggo nemtokake apa !!{formula} mawa kuantor
  universal~$\lforall[x][!A(x)]$ dicukupi, kita kudu nemtokake
  apa $\Sat{M}{\lnot\lexists[x][\lnot !A(x)]}[s]$.
  Upamane,
  \[
  \Sat{M}{\lforall[x][(R(x,a) \lif R(a,x))]}[s],
  \]
  Kaping pisan, $\Sat{M}{R(x,a) \lif R(a,x)}[\Subst{s}{m}{x}]$
  kanggo kabeh $m \in \Domain M$
  ($\Sat{M}{R(a,x)}[\Subst{s}{1}{x}]$ lan
  $\Sat/{M}{R(x,a)}[\Subst{s}{m}{x}]$ kanggo $m = 2$,
  $3$, utawa~$4$). Cathetan koreksi sumber OLPL-078:
  predikat ing alesan kanggo telung nilai pungkasan
  kudu duwe variabel ing argumen kapisan.
  Mula ora ana $m \in \Domain M$
  kang ndadekake $\Sat{M}{\lnot(R(x,a) \lif R(a,x))}[\Subst{s}{m}{x}]$,
  mula $\Sat/{M}{\lexists[x][\lnot(R(x,a) \lif R(a,x))]}[s]$.
  Dadi $\Sat{M}{\lnot\lexists[x][\lnot(R(x,a) \lif R(a,x))]}[s]$.
  Ing sisih liya,
  \[
  \Sat/{M}{\lforall[x][(R(a,x) \lif R(x,a))]}[s],
  \]
  amarga $\Sat/{M}{R(a,x) \lif R(x,a)}[\Subst{s}{2}{x}]$,
  mula $\Sat{M}{\lexists[x][\lnot(R(a,x) \lif R(x,a))]}[s]$.
  Kaya kang bisa dikira saka tuladha iki,
  $\Sat{M}{\lforall[x][!A(x)]}[s]$ yen lan mung yen
  $\Sat{M}{!A(x)}[\Subst{s}{m}{x}]$ kanggo saben
  $m \in \Domain M$.}{}

Kanggo kasus kang luwih rumit, gatekna
\[
\lforall[x][(R(a,x) \lif \lexists[y][R(x,y)])].
\]
Amarga $\Sat/{M}{R(a,x)}[\Subst{s}{3}{x}]$ lan
$\Sat/{M}{R(a,x)}[\Subst{s}{4}{x}]$, kasus kang kudu
nggatekake konsekuen kondisional mung $m = 1$
lan $m = 2$. Cathetan koreksi sumber OLPL-075:
sumber ngilangake huruf variabel ing nilai kapindho;
indeks pangaji kang padha dibalekake. Apa
$\Sat{M}{\lexists[y][R(x,y)]}[\Subst{s}{1}{x}]$
dicukupi? Iki bener yen ana paling ora siji $n \in \Domain M$
kang ndadekake $\Sat{M}{R(x,y)}[\Subst{\Subst{s}{1}{x}}{n}{y}]$.
Nyatane, yen kita njupuk $n = 1$, kita nduweni
$\Subst{\Subst{s}{1}{x}}{n}{y} = \Subst{s}{1}{y} =
s$. Amarga $s(x) = 1$, $s(y) = 1$, lan
$\tuple{1,1} \in \Assign{R}{M}$, wangsulane ya.

Kanggo nemtokake apa
$\Sat{M}{\lexists[y][R(x,y)]}[\Subst{s}{2}{x}]$,
kita kudu mriksa pangaji !!{variable}
$\Subst{\Subst{s}{2}{x}}{n}{y}$. Ing kene, kanggo $n = 1$,
pangaji iki yaiku~$s_2 = \Subst{s}{2}{x}$,
kang ora nyukupi $R(x,y)$ ($s_2(x) = 2$, $s_2(y) = 1$,
lan $\tuple{2,1}\notin \Assign{R}{M}$).
Nanging, gatekna $\Subst{\Subst{s}{2}{x}}{3}{y} =
\Subst{s_2}{3}{y}$.
$\Sat{M}{R(x,y)}[\Subst{s_2}{3}{y}]$ amarga
$\tuple{2,3} \in \Assign{R}{M}$, mula
$\Sat{M}{\lexists[y][R(x,y)]}[s_2]$.

Dadi, kanggo kabeh $m \in \Domain M$, salah siji
$\Sat/{M}{R(a,x)}[\Subst{s}{m}{x}]$ (yen $m = 3$, $4$)
utawa $\Sat{M}{\lexists[y][R(x,y)]}[\Subst{s}{m}{x}]$
(yen $m = 1$, $2$), mula
\[
\Sat{M}{\lforall[x][(R(a,x) \lif \lexists[y][R(x,y)])]}[s].
\]
Cathetan koreksi sumber OLPL-076: kuantifikasi panutup
sumber nggunakake huruf kanggo pangaji variabel njero,
dene saben kasus ing ukara iki miturut pangaji variabel njaba;
huruf kang cocog dibalekake. Ing sisih liya,
\[
\Sat/{M}{\lexists[x][(R(a,x) \land \lforall[y][R(x,y)])]}[s].
\]
Kita nduweni $\Sat{M}{R(a,x)}[\Subst{s}{m}{x}]$
mung kanggo $m = 1$ lan $m = 2$. Nanging kanggo
rong nilai~$m$ iku, ana maneh $n \in \Domain M$,
yaiku $n = 4$, kang ndadekake
$\Sat/{M}{R(x,y)}[\Subst{\Subst{s}{m}{x}}{n}{y}]$,
mula $\Sat/{M}{\lforall[y][R(x,y)]}[\Subst{s}{m}{x}]$
kanggo $m = 1$ lan $m = 2$. Dadi ora ana
$m \in \Domain M$ kang ndadekake
$\Sat{M}{R(a,x) \land \lforall[y][R(x,y)]}[\Subst{s}{m}{x}]$.
\end{ex}

\iftag{defEx}{%
\begin{prop}\ollabel{prop:sat-ex}
$\Sat{M}{\lexists[x][!B(x)]}[s]$ yen lan mung yen ana varian-$x$ $s'$ saka $s$
    kang ndadekake $\Sat{M}{!B(x)}[s']$.
\end{prop}

\begin{proof}
  Gladhen.
\end{proof}
}{}

\tagprob{defEx}
\begin{prob}
  Buktèkna \olref[fol][syn][sat]{prop:sat-ex}
\end{prob}
\tagendprob

\iftag{defAll}{%
\begin{prop}\ollabel{prop:sat-all}
$\Sat{M}{\lforall[x][!B(x)]}[s]$ yen lan mung yen kanggo saben varian-$x$~$s'$ saka $s$,
  $\Sat{M}{!B(x)}[s']$
\end{prop}

\begin{proof}
  Gladhen.
\end{proof}
}{}

\tagprob{defAll}
\begin{prob}
  Buktèkna \olref[fol][syn][sat]{prop:sat-all}
\end{prob}
\tagendprob

\begin{prob}
Upamakna $\Lang L = \{c, f, A\}$ kanthi siji !!{constant}, siji
!!{function} siji panggonan, lan siji !!{predicate} rong panggonan,
lan !!{structure}~$\Struct{M}$ ditemtokake dening
\begin{enumerate}
\item $\Domain M = \{1, 2, 3\}$
\item $\Assign{c}{M} = 3$
\item $\Assign{f}{M}(1) = 2, \Assign{f}{M}(2) = 3, \Assign{f}{M}(3) = 2$
\item $\Assign{A}{M} = \{\tuple{1, 2}, \tuple{2, 3}, \tuple{3, 3}\}$
\end{enumerate}
(a) Upamakna $s(v) = 1$ kanggo kabeh !!{variable}s~$v$.
Temtokna apa
\[
\Sat{M}{\lexists[x][(A(f(z), c) \lif \lforall[y][(A(y, x) \lor A(f(y),
      x))])]}[s]
\]
Jlentrehna alesane.

(b) Wenehana struktur lan pangaji !!{variable} liyane kang
!!{formula} iki ora dicukupi.
\end{prob}

\end{document}
